\section{Duplicate e Intervention: descomponer la learned lookup en pasos observables}

La corruption sobre datos reales muestra que el canal relacional se utiliza, pero aún no deja claro qué algoritmo aprendió el modelo. La semi-synthetic duplicate task construye un mecanismo más transparente: se copian varias rows, se ocultan los target de las originals y se exponen los target de las duplicates; una predicción correcta debe emparejar la row correspondiente, leer el target y luego copiarlo de vuelta a la query.

\lecturefigure{slide-27-duplicate-intervention.jpg}{Cadena completa de evidencia de la semi-synthetic duplicate task, la attention, las predicciones y la target intervention.}{Video oficial de Stanford Online, 00:54:44--01:00:41.}

\begin{knowledgebox}{Lectura de la figura: orden de razonamiento de los cinco panel}
\begin{enumerate}
  \item (a) Construir las original/duplicate rows y ocultar los targets de las originals;
  \item (b) verificar si la ABD attention se concentra en las duplicates correspondientes;
  \item (c) comprobar que las predictions tienen una correlación alta con los duplicate targets;
  \item (d) modificar deliberadamente un duplicate target en test time, sin volver a entrenar;
  \item (e) observar si la original prediction cambia de acuerdo con la intervention.
\end{enumerate}
\end{knowledgebox}

Esta tarea exige que el modelo aprenda un programa compuesto:
\[
x_i\xrightarrow{\text{match features}}
x_{j(i)}^{\mathrm{dup}}
\xrightarrow{\text{read target}}
y_{j(i)}^{\mathrm{dup}}
\xrightarrow{\text{copy/transform}}
\widehat y_i.
\]

\begin{itemize}
  \item $j(i)$: el duplicate index que corresponde a la original row $i$;
  \item el match lo realiza la learned attention representation, no un row ID escrito de antemano;
  \item la intervention modifica $y_{j(i)}^{\mathrm{dup}}$ y deja sin cambios el resto del mecanismo;
  \item si $\widehat y_i$ se desplaza de manera estable con la intervention, el modelo usa una lookup path generalizable.
\end{itemize}

\teachervoice{El ponente reporta que, después de la intervention, la correlación entre la prediction y el duplicate value se mantiene en aproximadamente 99.6\%. Lo decisivo es que el modelo, sin volver a entrenarse, sigue un target value que está fuera de la distribución de entrenamiento; esto se acerca más a una prueba del mecanismo que observar solo el test error habitual.}

\subsection{Por qué no es solo nearest neighbor}

Un plain k-NN también puede encontrar la duplicate y copiar el target, por lo que la tarea básica no basta para mostrar la transformación relacional aprendible de NPT. Los autores van más allá y le suman uno al duplicate target; una copy rule fija se equivoca de manera sistemática, mientras que NPT puede aprender a hacer primero la lookup y después restar uno. De forma más general, NPT puede aprender una joint feature--target relation, en lugar de solo promediar localmente bajo una metric fija.

\begin{importantbox}{Lo que aprende NPT no es una tabla estática de vecinos}
La evidencia más fuerte no es que «la attention se dirige a rows similares», sino que el modelo, después del match, todavía puede ejecutar una task-dependent transformation. Las múltiples capas ABD/ABA permiten que la lookup, el attribute reasoning y la value transformation se combinen en un learned algorithm.
\end{importantbox}

\teachervoice{Kossen reconoce por iniciativa propia que la duplicate task básica puede resolverse con nearest neighbor, y luego propone la transformación de sumar uno para aumentar la dificultad. Vale la pena aprender esta forma de argumentar: no evitar los baselines sencillos, sino construir el experimento mínimo que distinga la capacidad del mecanismo.}

\begin{warningbox}{Límites para extrapolar el semi-synthetic success}
La tarea se diseñó deliberadamente para que existan duplicates claras e interventions verificables, por lo que sirve para demostrar la representational capability. Los datos reales normalmente no tienen coincidencias perfectas, un mecanismo único y limpio ni labels manipulables; el éxito no equivale a que NPT ya domine el causal reasoning en general.
\end{warningbox}

\subsection{Resumen de la sección}

El duplicate/intervention experiment descompone el razonamiento entre datapoints en cuatro pasos: match, read, transform y write, y usa la test-time intervention para comprobar que la prediction cambia con el context. Es una evidencia fuerte del mecanismo, pero sigue limitada a una semi-synthetic construction transparente.
